% ============================================================================
%  অধ্যায় ৩ — উৎপাদক বিশ্লেষণ
%  COMPILE WITH XeLaTeX  (twice, so the point total resolves)
%  Overleaf: Menu -> Compiler -> XeLaTeX
% ============================================================================
\documentclass[11pt]{article}
\usepackage{bnhandout}          % add [nosol] for a student edition

\hauthor[https://sites.google.com/view/jugantordey/home]{Jugantor Dey}
\hseries{\textit{ Mathematics for the Physics Olympiad}}

\begin{document}

\handouttitle{Chapter 3 : Factorization}

\noindent
এই অধ্যায়ের পূর্বশর্ত অধ্যায় ১ ও ২। ২ থেকে তিনটে জিনিস এখানে প্রতি
পৃষ্ঠায় লাগবে : বণ্টন বিধি, অভেদ ও সমীকরণের পার্থক্য, আর পূর্ণসংখ্যা সূচকের
নিয়মগুলো। ১ থেকে লাগবে কেবল একটা অভ্যাস : দাবি করলে প্রমাণ করতে হয়।

অধ্যায় ২-এ আমরা বারবার একটা কাজ করেছি --- বন্ধনী খুলেছি। $(x+3)(x+5)$ থেকে
$x^2 + 8x + 15$। এই অধ্যায় পুরোটাই সেই কাজটা উল্টো দিকে করার গল্প: $x^2 + 8x
+ 15$ হাতে পেয়ে $(x+3)(x+5)$ ফিরে পাওয়া। শুনতে নিছক কৌশল মনে হয়, কিন্তু এর
পেছনে একটা গভীর কারণ আছে, আর সেটা এখনই বলে রাখা ভালো। গুণফল শূন্য হওয়ার একটাই
উপায় : কোনো একটা উৎপাদক শূন্য হওয়া। তাই একটা রাশিকে উৎপাদকে ভাঙতে পারা মানে
তার শূন্য হওয়ার সব শর্ত একসাথে জেনে ফেলা। পুরো সিরিজে সমীকরণ সমাধানের মূল
হাতিয়ার এটাই। এখানে মোট \totalpoints\ পয়েন্ট আছে।

\section{সাধারণ উৎপাদক}

\begin{idea}[উৎপাদক বিশ্লেষণ মানে গুণকে উল্টো পথে হাঁটানো]
কোনো রাশিকে উৎপাদকে বিশ্লেষণ করা মানে তাকে কতগুলো সরল রাশির \emph{গুণফল} হিসেবে
লেখা, এমনভাবে যেন নতুন লেখাটা পুরোনোটার সাথে একটা অভেদ হয়; অর্থাৎ চলকের সব
মানের জন্যই দুই পাশ সমান থাকে।

এই কারণেই উৎপাদক বিশ্লেষণ যাচাই করা সহজ। উত্তর ঠিক হলো কি না জানতে গুণ করে
দেখলেই হয়। ভুল হলে ধরা পড়বে, কারণ অভেদ ভাঙার জন্য একটা মানই যথেষ্ট।
\end{idea}

\begin{idea}[সাধারণ উৎপাদক বণ্টন বিধিরই উল্টো পিঠ]
অধ্যায় ২-এ আমরা ধরে নিয়েছিলাম $a(b+c) = ab + ac$। ডান থেকে বাঁয়ে পড়লে
এটাই সাধারণ উৎপাদক বের করার নিয়ম:
\[
  ab + ac = a(b + c).
\]
নতুন কোনো নিয়ম শেখার নেই। যে অংশটা প্রতিটি পদে আছে, সেটাকে বাইরে এনে বসাও।

একাধিক পদে সাধারণ অংশ খুঁজতে হলে দুটো জিনিস দেখো : সহগগুলোর গ.সা.গু., আর
প্রতিটি চলকের \emph{সবচেয়ে ছোট} সূচক। যেমন $12x^{4}y^{2} - 18x^{2}y^{3}$-এ
সহগের গ.সা.গু. $6$, $x$-এর ছোট সূচক $2$, $y$-এর ছোট সূচক $2$, তাই সাধারণ উৎপাদক
$6x^{2}y^{2}$।
\end{idea}

\noindent
কিন্তু প্রতিটি পদে সাধারণ কিছু না-ও থাকতে পারে, অথচ কিছু পদ আলাদা করে দলবদ্ধ
করলে থাকতে পারে। এই কৌশলের নাম দলবদ্ধকরণ (grouping)। ভাবনাটা সহজ: পদগুলোকে এমন
ভাগে ভাগ করো যেন প্রতি ভাগ থেকে সাধারণ উৎপাদক বের করার পর ভেতরে \emph{একই}
বন্ধনী পড়ে থাকে; তখন সেই বন্ধনীটাই গোটা রাশির সাধারণ উৎপাদক।

\begin{example}
উৎপাদকে বিশ্লেষণ করো: $\;2ax + 3ay - 2bx - 3by$।
\end{example}

\begin{solbox}
প্রথম দুটি পদে $a$ সাধারণ, শেষ দুটিতে $-b$:
\[
  2ax + 3ay - 2bx - 3by = a(2x + 3y) - b(2x + 3y).
\]
দুই জায়গায় একই বন্ধনী $(2x+3y)$ পড়ে আছে, তাই সেটাকেই বাইরে আনি:
\[
  = (2x + 3y)(a - b).
\]

শেষ দুই পদ থেকে $-b$ বের করাটা ইচ্ছাকৃত। $+b$ বের করলে পেতাম $b(-2x - 3y)$,
যেখানে বন্ধনীর অংশটা আগেরটার সাথে মেলে না। কোন চিহ্ন বাইরে আনব, সেটা ঠিক করে দেয়
লক্ষ্যটা : দুই বন্ধনী এক হওয়া চাই।

যাচাই: $(2x+3y)(a-b) = 2ax - 2bx + 3ay - 3by$, যা মূল রাশিই।
\end{solbox}

\begin{remark}
দলবদ্ধকরণে পদের ক্রম বদলানো লাগতে পারে, আর সেটা করা বৈধ। অধ্যায় ২-এ
আমরা $a + b = b + a$ ধরে নিয়েছি। কোন জোড়ায় ভাগ করবে তার কোনো সূত্র নেই; দু-একটা
সাজানো চেষ্টা করে দেখতে হয়। একটা কাজ না করলে অন্যটা করো, এটাকে ব্যর্থতা ভাবার
কিছু নেই।
\end{remark}

\prob{1} উৎপাদকে বিশ্লেষণ করো।
\begin{enumerate}
  \item $15a^{3}b^{2} - 25a^{2}b^{3} + 10a^{2}b^{2}$
  \item $px + py + qx + qy$
\end{enumerate}

\begin{sol}
\begin{enumerate}
  \item সহগের গ.সা.গু. $5$, $a$-এর ছোট সূচক $2$, $b$-এর ছোট সূচক $2$। তাই সাধারণ
        উৎপাদক $5a^{2}b^{2}$, এবং
        \[
          15a^{3}b^{2} - 25a^{2}b^{3} + 10a^{2}b^{2} = 5a^{2}b^{2}(3a - 5b + 2).
        \]
  \item $p(x+y) + q(x+y) = (x+y)(p+q)$।
\end{enumerate}
\end{sol}

\prob{2} উৎপাদকে বিশ্লেষণ করো: $\;x^{3} + 2x^{2} - 9x - 18$।

\begin{sol}
প্রথম দুই পদে $x^{2}$ আর শেষ দুই পদে $-9$ সাধারণ:
\[
  x^{2}(x + 2) - 9(x + 2) = (x + 2)(x^{2} - 9).
\]
এখানে থেমে যাওয়া যাবে না। $x^{2} - 9$ আরও ভাঙা যায়, কারণ $x^2 - 9 = x^2 - 3^2$।
পরের অনুচ্ছেদের অভেদ ব্যবহার করে
\[
  = (x + 2)(x - 3)(x + 3).
\]
নিয়ম হিসেবে মনে রাখো: প্রতিটি উৎপাদকের দিকে আবার তাকাও, আর কিছু ভাঙা যায় কি না
দেখো। ``উৎপাদকে বিশ্লেষণ করো'' মানে সবসময়ই \emph{সম্পূর্ণ} বিশ্লেষণ।
\end{sol}

\prob{3} উৎপাদকে বিশ্লেষণ করো: $\;x^{4} + x^{2} + 1$।

\begin{sol}
সরাসরি কোনো দলবদ্ধকরণ কাজ করে না। কৌশলটা হলো একটা পদ যোগ করে আবার বিয়োগ করা ---
এতে রাশির মান বদলায় না, কিন্তু চেহারা বদলায়। $x^{2}$ যোগ করে বিয়োগ করি:
\[
  x^{4} + x^{2} + 1 = (x^{4} + 2x^{2} + 1) - x^{2}.
\]
প্রথম বন্ধনীটা পুরো বর্গ, কারণ $x^{4} + 2x^{2} + 1 = (x^{2}+1)^{2}$। সুতরাং
\[
  = (x^{2}+1)^{2} - x^{2},
\]
যা দুই বর্গের অন্তর। তাই
\[
  = \big( (x^{2}+1) - x \big)\big( (x^{2}+1) + x \big)
  = (x^{2} - x + 1)(x^{2} + x + 1).
\]
যাচাই: $x = 1$ বসালে বাঁ পাশ $3$, ডান পাশ $1 \cdot 3 = 3$। ঠিক আছে।

কেন $x^{2}$ যোগ করলাম, আর কেন $2x^{2}$ বা $3x^{2}$ নয়? কারণ লক্ষ্য ছিল পূর্ণ বর্গ
বানানো। $x^{4} + 1$-কে বর্গ করতে ঠিক $2x^{2}$ দরকার, আর আমার হাতে ছিল $x^{2}$;
ঘাটতি $x^{2}$, সেটাই যোগ করে ধার শোধ করলাম।
\end{sol}

\section{মধ্যপদ বিশ্লেষণ}

\begin{idea}[পাঁচটি অভেদ, সবগুলোই বন্ধনী খুলে পাওয়া]
নিচের প্রতিটিই অভেদ, আর প্রতিটিরই প্রমাণ ডান পাশ খুলে ফেলা। কোনোটাই মুখস্থ করার
জিনিস নয়; ভুলে গেলে ত্রিশ সেকেন্ডে বের করে নেওয়া যায়।
\[
  (a+b)^{2} = a^{2} + 2ab + b^{2}, \qquad (a-b)^{2} = a^{2} - 2ab + b^{2},
\]
\[
  a^{2} - b^{2} = (a+b)(a-b), (a+b)^2-(a-b)^2=4ab
\]
\[
  a^{3} + b^{3} = (a+b)(a^{2} - ab + b^{2}), \qquad
  a^{3} - b^{3} = (a-b)(a^{2} + ab + b^{2}).
\]
শেষ দুটোর প্রমাণ দেখে নাও। $(a+b)(a^{2} - ab + b^{2})$ খুললে পাই
\[
  a^{3} - a^{2}b + ab^{2} + a^{2}b - ab^{2} + b^{3},
\]
যেখানে মাঝের চারটি পদ জোড়ায় জোড়ায় কেটে গিয়ে থাকে $a^{3} + b^{3}$। অন্যটাও একই
হিসাব, শুধু চিহ্ন বদলে।
\end{idea}

\begin{remark}
$a^{2} + b^{2}$ কিন্তু এই তালিকায় নেই, আর সেটা ভুলবশত নয়। বাস্তব সংখ্যার
মধ্যে ওটাকে ভাঙা যায় না। এই ``যায় না''-টা অধ্যায় ১৮-এ বদলে যাবে, যখন
আমরা $i$ চালু করব এবং $a^{2} + b^{2} = (a+ib)(a-ib)$ লিখতে পারব। উৎপাদক
বিশ্লেষণের উত্তর তাই শুধু প্রশ্নের উপরই নির্ভর করে না; কোন সংখ্যার জগতে কাজ করছি, তার উপরও নির্ভর করে।
\end{remark}

\begin{idea}[মধ্যপদ বিশ্লেষণ কেন কাজ করে]
$x^{2} + bx + c$ ভাঙতে চাইলে দেখো,
\[
  (x + p)(x + q) = x^{2} + (p + q)x + pq .
\]
তাই দরকার এমন দুটো সংখ্যা $p, q$ যাদের \emph{যোগফল} $b$ আর \emph{গুণফল} $c$।
পেয়ে গেলে উত্তর $(x+p)(x+q)$।

সহগ $a \neq 1$ হলে, অর্থাৎ $ax^{2} + bx + c$-এর ক্ষেত্রে, নিয়মটা একটু বদলায়:
এবার $p + q = b$ আর $pq = ac$। কেন? পুরো রাশিটাকে $a$ দিয়ে গুণ করে দেখো;
অধ্যায় ২-এর সূচক নিয়ম মেনে $a^{2}x^{2} = (ax)^{2}$, তাই
\[
  a(ax^{2} + bx + c) = (ax)^{2} + b(ax) + ac .
\]
ডান পাশে $ax$-কে একটা আস্ত রাশি ধরলে এটা আগের চেনা রূপ, যেখানে ``ধ্রুব পদ'' হলো
$ac$। তাই $p+q = b$, $pq = ac$ পেলে ডান পাশ দাঁড়ায় $(ax+p)(ax+q)$, এবং
\[
  ax^{2} + bx + c = \frac{(ax+p)(ax+q)}{a} .
\]
কার্যক্ষেত্রে ভগ্নাংশ টানাটানি না করে সহজ পথ হলো $bx$-কে $px + qx$ হিসেবে লিখে
দলবদ্ধকরণ প্রয়োগ করা --- ফল একই।
\end{idea}

\begin{example}
উৎপাদকে বিশ্লেষণ করো: $\;6x^{2} + 7x - 20$।
\end{example}

\begin{solbox}
এখানে $a = 6$, $b = 7$, $c = -20$, তাই $ac = -120$। দরকার দুটো সংখ্যা যাদের
গুণফল $-120$ আর যোগফল $7$। গুণফল ঋণাত্মক, তাই একটা ধনাত্মক একটা ঋণাত্মক; যোগফল
ধনাত্মক, তাই বড়টা ধনাত্মক। $120 = 15 \times 8$ আর $15 - 8 = 7$ --- পেয়ে গেছি:
$15$ ও $-8$।

মধ্যপদ ভেঙে দলবদ্ধ করি:
\[
  6x^{2} + 15x - 8x - 20 = 3x(2x + 5) - 4(2x + 5) = (2x + 5)(3x - 4).
\]
যাচাই: $(2x+5)(3x-4) = 6x^{2} - 8x + 15x - 20 = 6x^{2} + 7x - 20$। মিলেছে।

লক্ষ করো, $15x$ ও $-8x$-এর ক্রম উল্টে লিখলেও একই উত্তর আসত, শুধু বন্ধনী দুটো
অন্য ক্রমে বেরোত। ক্রম নিয়ে দুশ্চিন্তার কিছু নেই।
\end{solbox}

\prob{1} উৎপাদকে বিশ্লেষণ করো।
\begin{enumerate}
  \item $x^{2} + 9x + 20$
  \item $49a^{2} - 36b^{2}$
\end{enumerate}

\begin{sol}
\begin{enumerate}
  \item গুণফল $20$ ও যোগফল $9$ --- সংখ্যা দুটি $4$ ও $5$। তাই $(x+4)(x+5)$।
  \item $49a^{2} = (7a)^{2}$ এবং $36b^{2} = (6b)^{2}$, তাই দুই বর্গের অন্তর:
        $(7a + 6b)(7a - 6b)$।
\end{enumerate}
\end{sol}

\prob{2} উৎপাদকে বিশ্লেষণ করো।
\begin{enumerate}
  \item $6x^{2} - 5x - 6$
  \item $27a^{3} - 8b^{3}$
\end{enumerate}

\begin{sol}
\begin{enumerate}
  \item $ac = -36$, যোগফল $-5$। সংখ্যা দুটি $-9$ ও $4$। তাই
        \[
          6x^{2} - 9x + 4x - 6 = 3x(2x - 3) + 2(2x - 3) = (2x - 3)(3x + 2).
        \]
  \item $27a^{3} = (3a)^{3}$ ও $8b^{3} = (2b)^{3}$। ঘনের অন্তরের অভেদে
        $a \to 3a$, $b \to 2b$ বসিয়ে
        \[
          (3a - 2b)\big( (3a)^{2} + (3a)(2b) + (2b)^{2} \big)
          = (3a - 2b)(9a^{2} + 6ab + 4b^{2}).
        \]
\end{enumerate}
\end{sol}

\prob{2} $x^{4} - 16$ রাশিটিকে বাস্তব সংখ্যার মধ্যে সম্পূর্ণরূপে উৎপাদকে
বিশ্লেষণ করো, এবং ব্যাখ্যা করো কেন আর এগোনো যায় না।

\begin{sol}
$x^{4} = (x^{2})^{2}$ ও $16 = 4^{2}$, তাই প্রথমে
\[
  x^{4} - 16 = (x^{2} - 4)(x^{2} + 4).
\]
প্রথম উৎপাদকটি আবার দুই বর্গের অন্তর: $x^{2} - 4 = (x-2)(x+2)$। সুতরাং
\[
  x^{4} - 16 = (x - 2)(x + 2)(x^{2} + 4).
\]
$x^{2} + 4$ আর ভাঙা যায় না, কারণ যেকোনো বাস্তব $x$-এর জন্য $x^{2} \geq 0$,
তাই $x^{2} + 4 \geq 4 > 0$ --- রাশিটি কখনোই শূন্য হয় না। কিন্তু বাস্তব সহগে
$x^{2} + 4 = (x - r)(x - s)$ লেখা গেলে $x = r$ বসিয়ে বাঁ পাশ শূন্য পেতাম।
সুতরাং তেমন $r, s$ থাকতে পারে না।
\end{sol}

\prob{3} প্রমাণ করো যে সব বাস্তব $a, b, c$-এর জন্য
\[
  a^{3} + b^{3} + c^{3} - 3abc = (a + b + c)(a^{2} + b^{2} + c^{2} - ab - bc - ca).
\]
তারপর এই অভেদ ব্যবহার করে দেখাও: $a + b + c = 0$ হলে $a^{3} + b^{3} + c^{3} = 3abc$।

\begin{sol}
ডান পাশ খুলি। $(a+b+c)$-এর প্রতিটি পদকে বন্ধনীর ভেতরের প্রতিটি পদ দিয়ে গুণ করে
তিন সারিতে সাজাই:
\[
  a(a^{2} + b^{2} + c^{2} - ab - bc - ca)
  = a^{3} + ab^{2} + ac^{2} - a^{2}b - abc - a^{2}c,
\]
\[
  b(a^{2} + b^{2} + c^{2} - ab - bc - ca)
  = a^{2}b + b^{3} + bc^{2} - ab^{2} - b^{2}c - abc,
\]
\[
  c(a^{2} + b^{2} + c^{2} - ab - bc - ca)
  = a^{2}c + b^{2}c + c^{3} - abc - bc^{2} - ac^{2}.
\]
এবার তিন সারি যোগ করি। $ab^{2}$ প্রথম সারিতে ধনাত্মক, দ্বিতীয়টিতে ঋণাত্মক ---
কেটে গেল। একইভাবে $a^{2}b$, $ac^{2}$, $a^{2}c$, $bc^{2}$, $b^{2}c$ প্রতিটিই
জোড়ায় কেটে যায়। টিকে থাকে $a^{3} + b^{3} + c^{3}$, আর তিনটি সারি থেকে তিনটি
$-abc$, অর্থাৎ $-3abc$। তাই ডান পাশ $= a^{3} + b^{3} + c^{3} - 3abc$, যা বাঁ পাশ।

এখন $a + b + c = 0$ হলে ডান পাশের প্রথম বন্ধনী শূন্য, তাই গোটা গুণফল শূন্য:
\[
  a^{3} + b^{3} + c^{3} - 3abc = 0 \iff a^{3} + b^{3} + c^{3} = 3abc .
\]
\end{sol}

\section{উৎপাদক ও ভাগশেষ উপপাদ্য}

\noindent
এতক্ষণ আমরা চেনা ছাঁচ খুঁজেছি। কিন্তু $x^{3} - 6x^{2} + 11x - 6$ জাতীয় রাশির
সামনে পড়লে কোনো ছাঁচ কাজে আসে না। দরকার একটা পদ্ধতি, যা ছাঁচ চেনার উপর নির্ভর
করে না। সেটাই এই অনুচ্ছেদ।

\begin{idea}[$x^{k} - a^{k}$-এর সবসময় একটি উৎপাদক $(x - a)$]
যেকোনো ধনাত্মক পূর্ণসংখ্যা $k$-এর জন্য
\[
  x^{k} - a^{k} = (x - a)\big( x^{k-1} + x^{k-2}a + x^{k-3}a^{2} + \cdots + a^{k-1} \big).
\]
প্রমাণ: ডান পাশ খুলে দুই ভাগে ভাবো। $x$ দিয়ে গুণ করলে পাই
$x^{k} + x^{k-1}a + \cdots + xa^{k-1}$, আর $-a$ দিয়ে গুণ করলে পাই
$-x^{k-1}a - x^{k-2}a^{2} - \cdots - a^{k}$। দ্বিতীয় তালিকার প্রতিটি পদ প্রথম
তালিকার ঠিক পরের পদটিকে কেটে দেয়। বাকি থাকে কেবল প্রথম তালিকার প্রথম পদ $x^{k}$
আর দ্বিতীয় তালিকার শেষ পদ $-a^{k}$। ($k=2$ ও $k=3$-এ এটা আগের অনুচ্ছেদের দুই
বর্গের অন্তর ও ঘনের অন্তর।)
\end{idea}

\begin{idea}[ভাগশেষ উপপাদ্য]
যেকোনো বহুপদবিশিষ্ট রাশি বা বহুপদী (polynomial) $f(x)$ -কে $(x-a)$ দ্বারা ভাগ করলে ভাগশেষ $f(a)$।

প্রমাণ: মনে করি, $f(x)$-কে $(x-a)$ দ্বারা ভাগ করলে ভাগফল $q(x)$ এবং ভাগশেষ $r$, যেখানে $f(x), q(x)$ $x$-এর একেকটি বহুপদী। তাহলে,
\[
  f(x)=(x-a)q(x)+r.
\]
উপরের সমীকরণটিতে $x=a$ বসালে,
\[
  f(a)=(a-a)q(a)+r=0+r=r
\]
অর্থাৎ, ভাগশেষ $f(a)$।
\end{idea}

\begin{idea}[উৎপাদক উপপাদ্য]
আগেরটার সরাসরি ফল:
\[
  (x - a) \text{ হলো } f(x)\text{-এর একটি উৎপাদক} \iff f(a) = 0 .
\]
ডান থেকে বাঁ: $f(a) = 0$ হলে ভাগশেষ উপপাদ্য বলছে $f(x) = (x-a)q(x)$, অর্থাৎ
$(x-a)$ উৎপাদক। বাঁ থেকে ডান: $f(x) = (x-a)q(x)$ হলে $x = a$ বসিয়ে পাই
$f(a) = 0 \cdot q(a) = 0$।

লক্ষ করো এটা $\Leftrightarrow$, কেবল $\Rightarrow$ নয় : দুই দিকই আলাদা করে
দেখাতে হয়েছে, আর দুটোই লাগবে।
\end{idea}

\begin{remark}
কোন $a$ চেষ্টা করব? বহুপদীর সকল সহগ পূর্ণসংখ্যা আর সর্বোচ্চ ঘাতের সহগ $1$ হলে, $a$ এর মান যদি কোনো
পূর্ণসংখ্যাহয় তাহলে সেটা ধ্রুব পদের একটি উৎপাদক হতেই হবে। কারণটা দেখা কঠিন নয়:
$f(a) = 0$ লিখে ধ্রুব পদটা এক পাশে সরালে বাকি সব পদেই অন্তত একটা $a$ থাকে, তাই
$a$ ধ্রুব পদের একটি উৎপাদক হবে। এই ছোট পর্যবেক্ষণটাই $a$ এর জন্য মান খোঁজার তালিকা কয়েকটা সংখ্যায় নামিয়ে আনে।
 \end{remark}

\begin{example}
$f(x) = x^{3} - 7x + 6$ হলে $f(x)$-কে $(x - 2)$ দিয়ে ভাগ করলে ভাগশেষ কত?
তারপর $f(x)$-কে সম্পূর্ণরূপে উৎপাদকে বিশ্লেষণ করো।
\end{example}

\begin{solbox}
ভাগশেষ উপপাদ্য অনুযায়ী ভাগশেষ $= f(2) = 8 - 14 + 6 = 0$। মানে $(x-2)$
আসলে একটি উৎপাদক।

বাকি অংশ বের করতে $x^{3} - 7x + 6$ থেকে $(x-2)$ টেনে বের করি grouping দিয়ে।
$x^{3} - 7x + 6 = x^{3} - 2x^{2} + 2x^{2} - 4x - 3x + 6$, অর্থাৎ
\[
  x^{2}(x - 2) + 2x(x - 2) - 3(x - 2) = (x - 2)(x^{2} + 2x - 3).
\]
এবার মধ্যপদ বিশ্লেষণ: গুণফল $-3$, যোগফল $2$ --- সংখ্যা দুটি $3$ ও $-1$। তাই
$x^{2} + 2x - 3 = (x + 3)(x - 1)$, এবং
\[
  f(x) = (x - 2)(x + 3)(x - 1).
\]
যাচাই: $f(1) = 1 - 7 + 6 = 0$ ও $f(-3) = -27 + 21 + 6 = 0$। তিনটি মূলই মিলেছে।

$2$ দিয়ে শুরু করলাম কেন? ধ্রুব পদ $6$-এর ভাজকগুলো $\pm 1, \pm 2, \pm 3, \pm 6$
--- চেষ্টা করার মতো সংখ্যা মোট আটটি, তার বেশি নয়।
\end{solbox}

\prob{2} $f(x) = 2x^{3} - 5x^{2} + 3x - 7$ হলে
\begin{enumerate}
  \item $(x - 2)$ দিয়ে ভাগ করলে ভাগশেষ কত?
  \item $(x + 1)$ কি $f(x)$-এর একটি উৎপাদক? কারণ দেখাও।
\end{enumerate}

\begin{sol}
\begin{enumerate}
  \item ভাগশেষ $= f(2) = 2(8) - 5(4) + 3(2) - 7 = 16 - 20 + 6 - 7 = -5$।
  \item $(x+1) = (x - (-1))$, তাই দেখতে হবে $f(-1)$ শূন্য কি না:
        \[
          f(-1) = 2(-1) - 5(1) + 3(-1) - 7 = -2 - 5 - 3 - 7 = -17 .
        \]
        শূন্য নয়, তাই উৎপাদক উপপাদ্য অনুযায়ী $(x+1)$ উৎপাদক নয়।
\end{enumerate}
\end{sol}

\prob{2} $k$-এর কোন মানের জন্য $(x - 3)$ রাশিটি
$\;x^{3} - 4x^{2} + kx - 6$-এর একটি উৎপাদক হবে?

\begin{sol}
$g(x) = x^{3} - 4x^{2} + kx - 6$ ধরি। উৎপাদক উপপাদ্য অনুযায়ী $(x-3)$ উৎপাদক
হবে তখনই, যখন $g(3) = 0$:
\[
  g(3) = 27 - 36 + 3k - 6 = 3k - 15 .
\]
সুতরাং $3k - 15 = 0$, অর্থাৎ $k = 5$।

যাচাই: $k = 5$ হলে $g(3) = 27 - 36 + 15 - 6 = 0$। ঠিক আছে।
\end{sol}

\prob{3} $\;x^{3} - 6x^{2} + 11x - 6\;$ রাশিটিকে সম্পূর্ণরূপে উৎপাদকে বিশ্লেষণ
করো।

\begin{sol}
$h(x) = x^{3} - 6x^{2} + 11x - 6$ ধরি। ধ্রুব পদ $-6$, তাই পূর্ণসংখ্যা মূল থাকলে
তা $\pm 1, \pm 2, \pm 3, \pm 6$-এর মধ্যে।
\[
  h(1) = 1 - 6 + 11 - 6 = 0 .
\]
প্রথম চেষ্টাতেই পাওয়া গেল, তাই $(x - 1)$ একটি উৎপাদক। বাকিটা বের করতে
দলবদ্ধকরণ:
\[
  x^{3} - x^{2} - 5x^{2} + 5x + 6x - 6
  = x^{2}(x-1) - 5x(x-1) + 6(x-1)
  = (x-1)(x^{2} - 5x + 6).
\]
এখন মধ্যপদ বিশ্লেষণ: গুণফল $6$, যোগফল $-5$ --- সংখ্যা দুটি $-2$ ও $-3$। সুতরাং
\[
  h(x) = (x - 1)(x - 2)(x - 3).
\]
যাচাই: $h(2) = 8 - 24 + 22 - 6 = 0$ এবং $h(3) = 27 - 54 + 33 - 6 = 0$।
\end{sol}

\section{দ্বিঘাত সমীকরণ}

\noindent
এতক্ষণ যা করলাম তার সবটুকুই ছাঁচ চেনার উপর নির্ভরশীল ছিল। মধ্যপদ বিশ্লেষণে
দুটো সুন্দর সংখ্যা খুঁজে পেতে হয়। কিন্তু $x^{2} - 4x + 1$-এর জন্য গুণফল $1$ আর
যোগফল $-4$ দেওয়া কোনো পূর্ণসংখ্যা জোড়া নেই। রাশিটার কি তবে মূল নেই? আছে, শুধু
সেগুলো পূর্ণসংখ্যা নয়। দরকার এমন একটা পদ্ধতি যা আন্দাজের উপর নির্ভর করে না।

\begin{idea}[পূর্ণবর্গ বানানো]
যেকোনো $x^{2} + px$-কে একটা পুরো বর্গের অংশ বানানো যায়, কারণ
$(x + \tfrac{p}{2})^{2} = x^{2} + px + \tfrac{p^{2}}{4}$। ঘাটতিটুকু যোগ করে
আবার বিয়োগ করে নিলেই হলো:
\[
  x^{2} + px = \left( x + \frac{p}{2} \right)^{2} - \frac{p^{2}}{4}.
\]
এর নাম পূর্ণবর্গ সম্পূর্ণ করা (completing the square)। লাভটা হলো, ডান পাশে
$x$ একবারমাত্র আছে বর্গের ভেতরে। ফলে সমীকরণ সমাধান করতে আর ছাঁচ চেনা লাগে
না, শুধু বর্গমূল নিলেই চলে।
\end{idea}

\begin{remark}[যা এখানে ধরে নেওয়া হচ্ছে]
প্রতিটি অঋণাত্মক বাস্তব সংখ্যা $D$-এর জন্য এমন একটিমাত্র অঋণাত্মক বাস্তব সংখ্যা
আছে যার বর্গ $D$ : একে লিখি $\sqrt{D}$। এই অস্তিত্বটা আমরা এখানে \emph{ধরে
নিচ্ছি}, প্রমাণ করছি না; প্রমাণের জন্য বাস্তব সংখ্যার পূর্ণতা নিয়ে কথা বলতে হয়,
যা এই সিরিজের অনেক পরের বিষয়। $D < 0$ হলে এমন কোনো বাস্তব সংখ্যা নেই, কারণ
যেকোনো বাস্তব সংখ্যার বর্গ অঋণাত্মক।
\end{remark}

\begin{idea}[দ্বিঘাত সমীকরণের সূত্র --- প্রমাণসহ]
$a \neq 0$ ধরে $ax^{2} + bx + c = 0$ সমাধান করি। প্রথমে দুই পাশকে $a$ দিয়ে ভাগ
করি (বৈধ, কারণ $a \neq 0$):
\[
  x^{2} + \frac{b}{a}x + \frac{c}{a} = 0 .
\]
পুরো বর্গ সম্পূর্ণ করি, $p = b/a$ নিয়ে:
\[
  \left( x + \frac{b}{2a} \right)^{2} - \frac{b^{2}}{4a^{2}} + \frac{c}{a} = 0,
  \qquad\text{অর্থাৎ}\qquad
  \left( x + \frac{b}{2a} \right)^{2} = \frac{b^{2} - 4ac}{4a^{2}} .
\]
ডান পাশের লবকে বলে নিশ্চায়ক (discriminant), $D = b^{2} - 4ac$। হর $4a^{2} > 0$,
তাই ডান পাশের চিহ্ন পুরোপুরি $D$-এর চিহ্নেই ঠিক হয়।

$D < 0$ হলে বাঁ পাশ একটা বর্গ, অর্থাৎ অঋণাত্মক, আর ডান পাশ ঋণাত্মক --- কোনো
বাস্তব $x$-এই সমতা সম্ভব নয়। বাস্তব সমাধান নেই।

$D \geq 0$ হলে বর্গমূল নিতে পারি। কোনো সংখ্যার বর্গ $t^{2}$ হলে সংখ্যাটা
$t$ বা $-t$, তাই
\[
  x + \frac{b}{2a} = \pm \frac{\sqrt{D}}{2a},
  \qquad\text{অর্থাৎ}\qquad
  x = \frac{-b \pm \sqrt{b^{2} - 4ac}}{2a} .
\]
($a$ ঋণাত্মক হলে $\sqrt{4a^{2}} = 2|a| = -2a$, কিন্তু সামনে $\pm$ থাকায় দুটো
উত্তরের \emph{জোড়া} একই থাকে --- শুধু কোনটা ধনাত্মক চিহ্ন থেকে এল সেটা বদলায়।
তাই $2a$ লিখলেই চলে।)
\end{idea}

\begin{idea}[নিশ্চায়ক কী বলে]
$D = b^{2} - 4ac$ মূল না বের করেই মূলের সংখ্যা বলে দেয়:
\begin{enumerate}
  \item $D > 0$ : দুটি ভিন্ন বাস্তব মূল।
  \item $D = 0$ : একটিমাত্র বাস্তব মূল, $x = -b/2a$; তখন
        $ax^{2}+bx+c = a(x + \tfrac{b}{2a})^{2}$, অর্থাৎ রাশিটি পুরো বর্গ।
  \item $D < 0$ : কোনো বাস্তব মূল নেই।
\end{enumerate}
তিনটেই উপরের প্রমাণ থেকে সরাসরি আসে, আলাদা করে মনে রাখার কিছু নেই।
\end{idea}

\begin{example}
পূর্ণবর্গ সম্পূর্ণ করে সমাধান করো: $\;x^{2} - 4x + 1 = 0$।
\end{example}

\begin{solbox}
এখানে $p = -4$, তাই $p/2 = -2$:
\[
  x^{2} - 4x + 1 = (x - 2)^{2} - 4 + 1 = (x - 2)^{2} - 3 .
\]
সমীকরণ দাঁড়ায় $(x-2)^{2} = 3$, তাই $x - 2 = \pm\sqrt{3}$, অর্থাৎ
\[
  x = 2 + \sqrt{3} \quad\text{অথবা}\quad x = 2 - \sqrt{3}.
\]
মধ্যপদ বিশ্লেষণ এখানে কেন খাটেনি তা এখন পরিষ্কার : মূল দুটি পূর্ণসংখ্যা নয়,
এমনকি ভগ্নাংশও নয়। ছাঁচ খোঁজার পদ্ধতি ব্যর্থ হয়েছিল পদ্ধতির দোষে নয়, উত্তরের
চেহারার কারণে।

সূত্র দিয়েও মিলিয়ে দেখা যায়: $D = 16 - 4 = 12$ এবং
$x = \dfrac{4 \pm \sqrt{12}}{2} = 2 \pm \sqrt{3}$, কারণ $\sqrt{12} = 2\sqrt{3}$।
\end{solbox}

\prob{2} পূর্ণবর্গ সম্পূর্ণ করে সমাধান করো: $\;2x^{2} - 7x + 3 = 0$। সূত্র
ব্যবহার কোরো না।

\begin{sol}
প্রথমে $2$ বাইরে আনি:
\[
  2\left( x^{2} - \frac{7}{2}x \right) + 3 = 0 .
\]
ভেতরে $p = -7/2$, তাই $p/2 = -7/4$ এবং
\[
  x^{2} - \frac{7}{2}x = \left( x - \frac{7}{4} \right)^{2} - \frac{49}{16}.
\]
বসিয়ে পাই
\[
  2\left( x - \frac{7}{4} \right)^{2} - \frac{49}{8} + 3 = 0,
  \qquad\text{অর্থাৎ}\qquad
  2\left( x - \frac{7}{4} \right)^{2} = \frac{25}{8}.
\]
দুই পাশকে $2$ দিয়ে ভাগ করে $\left( x - \dfrac{7}{4} \right)^{2} = \dfrac{25}{16}$,
তাই $x - \dfrac{7}{4} = \pm\dfrac{5}{4}$ এবং
\[
  x = \frac{7 + 5}{4} = 3 \qquad\text{অথবা}\qquad x = \frac{7 - 5}{4} = \frac{1}{2}.
\]
যাচাই: $2(9) - 7(3) + 3 = 0$ এবং $2\left(\tfrac{1}{4}\right) - \tfrac{7}{2} + 3 = 0$।
\end{sol}

\prob{3} $k$-এর কোন কোন মানের জন্য $\;kx^{2} + (k+2)x + 3 = 0\;$ সমীকরণটির ঠিক
একটি বাস্তব মূল থাকবে? সব ক্ষেত্র বিবেচনা করো।

\begin{sol}
দুটি আলাদা ক্ষেত্র, আর দ্বিতীয়টা বাদ দেওয়াই এখানে সবচেয়ে সাধারণ ভুল।

\textbf{ক্ষেত্র ১: $k = 0$।} তখন সমীকরণটা দ্বিঘাত নয়, রৈখিক: $2x + 3 = 0$,
যার একমাত্র মূল $x = -3/2$। ঠিক একটি বাস্তব মূল --- শর্ত পূরণ হলো। সুতরাং
$k = 0$ একটি উত্তর।

\textbf{ক্ষেত্র ২: $k \neq 0$।} এখন সমীকরণটা সত্যিই দ্বিঘাত, আর ঠিক একটি মূল
থাকার শর্ত $D = 0$:
\[
  D = (k+2)^{2} - 4 \cdot k \cdot 3 = k^{2} + 4k + 4 - 12k = k^{2} - 8k + 4 .
\]
$k^{2} - 8k + 4 = 0$ সমাধান করি পূর্ণবর্গ সম্পূর্ণ করে:
\[
  (k - 4)^{2} - 16 + 4 = 0 \implies (k-4)^{2} = 12 \implies k = 4 \pm 2\sqrt{3},
\]
কারণ $\sqrt{12} = 2\sqrt{3}$। দুটোই অশূন্য (আসন্ন মান $7.46$ ও $0.54$), তাই
দুটোই গ্রহণযোগ্য।

সুতরাং উত্তর: $k = 0$, $\;k = 4 + 2\sqrt{3}$, $\;k = 4 - 2\sqrt{3}$।

মনে রাখার কথাটা এই --- ``$D = 0$ দিলেই এক মূল'' নিয়মটা কেবল $a \neq 0$ ধরে
প্রমাণ করা হয়েছিল। যেখানে সর্বোচ্চ ঘাতের সহগে অক্ষর আছে, সেখানে সেই সহগ শূন্য
হওয়ার ক্ষেত্রটা আলাদা করে দেখতেই হবে।
\end{sol}

\begin{remark}
পূর্ণবর্গ সম্পূর্ণ করা এখানেই শেষ নয়। অধ্যায় ১২-এ দেখবে $x^{2} + px + q$-কে
বর্গের আকারে লেখা মানে আসলে parabola-র শীর্ষবিন্দু খুঁজে ফেলা, আর অধ্যায় ১৭-তে
দেখবে জটিল integral সরল করার প্রথম পদক্ষেপই প্রায়ই হরের পূর্ণবর্গ সম্পূর্ণ করা।
আর $D < 0$ হলে ``বাস্তব মূল নেই'' : এই দেয়ালটা অধ্যায় ১৮-তে ভাঙবে, যেখানে
দেখা যাবে মূল দুটি সবসময়ই আছে, শুধু বাস্তব সংখ্যার রেখার বাইরে।
\end{remark}

\end{document}